\uuid{N8uP}
\exo7id{1266}
\titre{exo7 1266}
\auteur{legall}
\organisation{exo7}
\datecreate{2003-10-01}
\isIndication{false}
\isCorrection{false}
\chapitre{Développement limité}
\sousChapitre{Applications}
\cpge{mpsi-mp2i.fonctions-valeurs-reelles-complexes.a, pcsi.techniques-calcul-fonctions.a, ptsi.techniques-calcul-fonctions.a, mpsi-mp2i.fonctions-valeurs-reelles-complexes.b, pcsi.techniques-calcul-fonctions.b, ptsi.techniques-calcul-fonctions.b, mpsi-mp2i.fonctions-valeurs-reelles-complexes.c, pcsi.techniques-calcul-fonctions.c, ptsi.techniques-calcul-fonctions.c, bcpst1.fonctions-reelles-usuelles.a, bcpst1.fonctions-reelles-usuelles.b, bcpst1.limites-continuite.c}
\module{Analyse}
\niveau{L1}
\difficulte{}

\contenu{
\begin{enumerate}
    \item \question{Soit $g$ la fonction d\'efinie par~: $$\displaystyle{ g(x) 
=\frac{x+1}{1+x^2} +\Arctan x}.$$
}
    \item \question{Quel est le domaine de d\'efinition de $g$~?}
    \item \question{Etudier ses
variations.}
    \item \question{Montrer que $g$ s'annule une et une seule fois sur $\Rr$ en un 
point $\alpha $ compris entre $-1$ et $0$ (on ne demande pas de 
pr\'eciser la valeur de $\alpha$).}
    \item \question{Dessiner le graphe de $g$.}
\end{enumerate}
}
